% 本文件由 LaTeX 排版 Agent 根据有效 translation.json 自动生成。
% author=Apocrypha; work_id=0811; title=交出般雀比辣多

\chapter{交出般雀比辣多}

% structure: p0001 source=h1 target=omitted reason=page-title-already-rendered-by-parent

那些文书送到了罗马城，呈递给凯撒，旁边站着不少人读完后，众人都惊愕不已，因为比辣多的邪恶导致黑暗和地震遍及全世界。凯撒满腔怒火，派兵丁去把比辣多押解前来。\par

比辣多被带到罗马城，凯撒听说他到了，就坐在众神庙中，全体元老院、全部军队和他的所有权贵都在场；凯撒命比辣多向前站。凯撒对他说：\Quote{你这极其不虔敬的人，既然在那个人身上看见了如此伟大的奇迹，怎敢做出这样的事？你胆敢行恶，已经毁灭了全世界。}\par

比辣多说：\Quote{全能的君王，这些事我并无罪过；但犹太人的群众是强暴有罪的。}凯撒说：\Quote{他们是谁？}比辣多说：\Quote{黑落德、阿尔赫劳、斐理伯、亚纳斯、盖法，以及犹太人的全体群众。}凯撒说：\Quote{你为何听从他们的建议？}比辣多说：\Quote{他们的民族叛逆不顺服，不服从您的权下。}凯撒说：\Quote{他们把他交给您时，您本应将他看管起来，送到我这里，而不应听从他们，把一个像他这样的义人、行了如此美好奇迹的人（正如您在报告中所说）钉在十字架上。因为这些奇迹显然表明耶稣就是基督，是犹太人的君王。}\par

正当凯撒这样说的时候，他刚一提到基督的名字，全体众神就一起倒坍，在凯撒与元老院所在之处化为尘土。站在凯撒身边的百姓都因这话语和众神的倾倒而战栗；他们恐惧万分，各自回家，对所发生的事感到惊奇。凯撒下令将比辣多严密看押，以便查明耶稣的真相。\par

第二天，凯撒坐在卡比托利欧山上，与全体元老院一同再次审问比辣多。凯撒说：\Quote{极其不虔敬的人，你要说实话，因为你向耶稣所行的不虔之举，甚至在这里也因众神的倾覆而显明了你的恶行。说吧，被钉死的那人是谁？因为连他的名字都毁灭了所有众神。}比辣多说：\Quote{关于他的记录确实是真的；因为我本人从他的作为中确信，他比我们所敬拜的众神都更伟大。}凯撒说：\Quote{那么，你为什么对他如此胆大妄为，做了这样的事？你难道不是明知他，却存心要危害我的王国吗？}比辣多说：\Quote{是因那些无法无天、不虔敬的犹太人的邪恶和叛乱，我才这样做的。}\par

凯撒满腔怒火，召集全体元老院和权贵商议，下令拟定一道针对犹太人的诏书如下：\Quote{致东境主要地区的总督利恰努斯，问安。眼下耶路撒冷居民及周围犹太城邑所行的鲁莽之事，以及他们的恶行，我已得知：他们强迫比辣多钉死了一位名叫耶稣的神祇；因他们的重大过犯，世界已变得黑暗，趋于毁灭。你要速速率领大队兵丁前往，按照这道诏书将他们擒拿。你要服从命令，采取行动对付他们，把他们驱散，使他们成为万民中的奴隶；将他们逐出整个犹太地区，使他们成为最微小的民族，以致再也看不见他们，因为他们充满了邪恶。}\par

这道诏书传到东方地区，利恰努斯因惧怕诏书而服从，擒拿了全体犹太民族；又将犹太地区剩余的人分散到各国，卖为奴隶。凯撒得知利恰努斯在东方地区对犹太人做了这些事，就高兴了。\par

凯撒又审问比辣多；他命令一位名叫阿尔比乌斯的队长砍下比辣多的头，说：\Quote{正如他下手对付那位名叫基督的义人，同样他也要这样跌倒，不得安全。}\par

比辣多前往刑场，默默祈祷说：\Quote{主，请不要把我与邪恶的希伯来人一同消灭，因为若不是那些无法无天的犹太人民族向我煽动叛乱，我本不会下手对付你。但你知道我是出于无知而行的。求你不要因这罪毁灭我；主，求你不记念我的恶，也不要记念你的仆人普洛克拉，她现在与我一同站在这死亡的时刻，你曾命她预言你将被钉在十字架上。求你不要因我的罪而判她受罚，但求你宽恕我们，使我们得以列在你义人的份中。}\par

看，比辣多祈祷完毕，就有声音从天上来，说：\Quote{万国的世代和家族都要称你有福，因为在你任内，先知论到我所说的一切话都应验了；你自己要在我第二次显现时成为我的见证人，那时我要审判以色列十二支派和那些不承认我名的人。}队长砍下了比辣多的头；看，主的天使接住了它。他的妻子普洛克拉看见天使来接他的头，自己也满怀喜乐，立即断了气，与她丈夫一同下葬。\par

\SourceRecord{Translated by Alexander Walker.；Ante-Nicene Fathers；Vol. 8.；Edited by Alexander Roberts, James Donaldson, and A. Cleveland Coxe.；Buffalo, NY: Christian Literature Publishing Co.,；1886.；<http://www.newadvent.org/fathers/0811.htm>.}
